% !PRINTED 506
\weberpage{506}

It has further been shown in \S\,134 that for $D\equiv 1\ (\mathrm{mod}\ 4)$
$\sqrt{D}\,\gamma_3(\omega)$ is a class invariant. But we have (\S\,54)

\weq{4}{\gamma_3(\omega) = \dfrac{8(2+\varkappa^2\varkappa'^2)(\varkappa'^2-\varkappa^2)}{\varkappa^2\varkappa'^2}.}

If $\omega$ satisfies equation (2), and one sets:

\weq{5}{\omega = \dfrac{\omega'-1}{\omega'+1},}

then for $\omega'$ the equation results:

\weq{6}{(A+B+C)\omega'^2 - 2(A-C)\omega' + (A-B+C) = 0.}

If $D\equiv 5\ (\mathrm{mod}\ 8)$, then $A$, $B$, $C$ are odd, and if
$D\equiv 1\ (\mathrm{mod}\ 8)$, we may assume $A$ and $C$ to be even, and
consequently in both cases (6) is a primitive equation for $\omega'$
of discriminant $4D$. Therefore, by \S\,126, 1., $f(\omega')^{24}$ is
a class invariant for this discriminant, and since now
by \S\,34 (18)

\[ f(\omega)f(\omega') = \sqrt{2} \]

holds, the same is true of $f(\omega)^{24}$. From this it follows that $\varkappa^2\varkappa'^2$
is a class invariant for $4D$, and since by (4) $\varkappa^2$ can be expressed rationally
in terms of $\gamma_3$ and $\varkappa^2\varkappa'^2$, it follows that:

3. \gesperrt{If the discriminant $D\equiv 1\ (\mathrm{mod}\ 4)$, then
$\varkappa^2$ is rationally expressible by a class invariant of the discriminant $4D$ and by $\sqrt{D}$.}

If $j(\omega)$ is a class invariant of the discriminant $D$, then
$j(2\omega)$, $j(\frac{1}{2}\omega)$ are class invariants for $4D$, and we have:

\weq{7}{
\begin{aligned}
j(2\omega) &= \dfrac{(f_1(2\omega)^{24}+16)^3}{f_1(2\omega)^{24}} = \dfrac{[256+f_2(\omega)^{24}]^3}{f_2(\omega)^{48}} \\
&= \dfrac{16(16\varkappa'^2+\varkappa^4)^3}{\varkappa^8\varkappa'^2}, \\
j\left(\dfrac{\omega}{2}\right) &= \dfrac{\left[f_2\!\left(\dfrac{\omega}{2}\right)^{24}+16\right]^3}{f_2\!\left(\dfrac{\omega}{2}\right)^{24}} = \dfrac{[256+f_1(\omega)^{24}]^3}{f_1(\omega)^{48}} \\
&= \dfrac{16(16\varkappa^2+\varkappa'^4)^3}{\varkappa^2\varkappa'^8}.
\end{aligned}
}

From this it follows that conversely the class invariants
of the discriminant $4D$ are also rationally expressible in terms of $\varkappa^2$.
